\documentclass[11pt]{article}
\usepackage[a4paper, top=18mm, bottom=18mm, left=20mm, right=20mm]{geometry}
\usepackage[T2A]{fontenc}
\usepackage[utf8]{inputenc}
\usepackage[ukrainian]{babel}
\usepackage{graphicx}
\usepackage[export]{adjustbox}
\usepackage{amsmath}
\usepackage{amssymb}
\usepackage{microtype}
\graphicspath{{/app/Quiz/Scripts/2023/ProgAll/15BinTree/}}
\setlength{\parindent}{0pt}
\setlength{\parskip}{6pt}
\emergencystretch=3em
\pagestyle{empty}
\begin{document}
%begin
{\centering\Large\bfseries Контрольна робота\par}
\medskip
\noindent\rule{\linewidth}{0.5pt}
\smallskip
\noindent\textbf{Варіант \No\,4}\hfill \textit{ПІБ:}\,\underline{\hspace{60mm}}
\vspace{4mm}

%beginitem
1. \textbf{Що буде виведено за виконання фрагменту коду?}

print('R', end=' ')
for b in range(14,14,-3):
    if b<=4:
        break
        print(b, end=' ')
    if b<=4:
        break
    else:
        print(b, end=' ')
    print(b)
\vspace{5mm}
%enditem

%beginitem
2. \textbf{Що буде виведено за виконання фрагменту коду?}
a = 1
b = 9
c = 4

def h(b):
global c    a *= 3
    b = 1
    c = 5
    return a + b + c

a = 5
b = 7
c = 3

try:
    print(h(a), a, b, c, sep=';')
except:
    print('Z', a, b, c, sep=';')
\vspace{5mm}
%enditem

%beginitem
3. 

\begin{center}
	\adjustimage{max size={0.8\linewidth}{0.8\paperheight}}
	{../15BinTree/trees_exp/52.png}
\end{center}
Указати послідовність міток вершин дерева, зображеного на рис., що утвориться в результаті обходу дерева в оберненому порядку. Мітки записувати через пробіл.\par Алгоритм починає роботу з кореня (на рис. це верхня вершина).
%answer:52 оберненому
\vspace{5mm}
%enditem

%beginitem
4. 
Записати ВИРАЗ, значенням якого є список, що складається з кубів цілих чисел від 12 до 2002 (включно), що не діляться на жодне з чисел 2, 3, записаних в оберненому порядку.\vspace{5mm}
%enditem

%beginitem
5. %goodpath:Scripts/2019/Mod2019_1/Lex/01-good.txt
\textbf{Відмітити все, що є однією правильною лексемою:}
( 1 ) true

( 2 ) 1+2j

( 3 ) xy

( 4 ) 1+1j
\vspace{5mm}
%enditem

%beginitem
6. 
Написати програму, що запитує у користувача значення дійсних параметрів $x$ та $y$, обчислює для них значення виразу 
$$\frac{\log_{2} x - 19}{(\tg x - 0.3) (y - 64)}$$ 
та повiдомляє введенi данi та результат обчислення.
 Оцінюється правильність та якість коду.

\vspace{5mm}
%enditem

%beginitem
7. \textbf{Що буде виведено за виконання фрагменту коду?}
b=['0','1','2','3','4']; print(b.pop(-4),end=' '); print(b)\vspace{5mm}
%enditem

%beginitem
8. 
Записати ВИРАЗ, значенням якого є множина, що складається з цілих чисел від 12 до 2002 (включно), що діляться на 3.\vspace{5mm}
%enditem

%beginitem
9. %goodpath:Scripts/2018/Mod2018_1/01-good.txt
\textbf{Відмітити все, що є однією правильною лексемою:}
( 1 ) %

( 2 ) x+23

( 3 ) True

( 4 ) 2int
\vspace{5mm}
%enditem

%beginitem
10. %goodpath:Scripts/2019/Mod2019_1/Lex/03-good.txt
\textbf{Відмітити все, що є коректним оператором С++:}
( 1 ) 3**-1

( 2 ) if ('x'><'y'): a=6

( 3 ) x<y<z

( 4 ) (2+2j)%2
\vspace{5mm}
%enditem

%end
\newpage
\end{document}

